% === Begin Label Data ===
% ID: 296
% 难度星级: 
% 标签: 
% 备注: 
% 组卷引用次数: 0
% === End  Label Data ===

\begin{problem}{2005}{G}{上海卷}{22}{圆锥曲线}
（1）求右焦点坐标是$(2,0)$，且经过点$(-2,-\sqrt{2})$的椭圆的标准方程；

（2）已知椭圆$C$的方程是$\dfrac{x^{2}}{a^{2}}+\dfrac{y^{2}}{b^{2}}=1(a>b>0)$．设斜率为$k$的直线$l$交椭圆$C$于$A$、$B$两点，$AB$的中点为$M$．证明：当直线$l$平行移动时，动点$M$在一条过原点的定直线上；

（3）利用（2）所揭示的椭圆几何性质，用作图方法找出下面给定椭圆的中心，简要写出作图步骤，并在图中标出椭圆的中心．

\begin{tikzpicture}[scale=1.2]
    % 1. 基础参数与椭圆
    \def\a{3.2}     % 半长轴
    \def\b{2.0}     % 半短轴
    \def\rot{30}    % 椭圆旋转角度
    
    % 绘制椭圆并命名路径以便求交点
    \draw[thick, rotate=\rot, name path=ellipse] (0,0) ellipse (\a cm and \b cm);
\end{tikzpicture}
\end{problem}

\begin{answer}
（1）$\dfrac{x^{2}}{8}+\dfrac{y^{2}}{4}=1$；（2）动点$M$在过原点的定直线$y=-\dfrac{b^{2}}{a^{2}k}x$上；（3）见解析．
\end{answer}

\begin{solutions}
（1）  
设椭圆的标准方程为$\dfrac{x^{2}}{a^{2}}+\dfrac{y^{2}}{b^{2}}=1(a>b>0)$，由右焦点坐标为$(2,0)$可得$c=2$，由椭圆的性质可得$a^{2}-b^{2}=c^{2}=4$．  

将点$(-2,-\sqrt{2})$代入椭圆方程得$\dfrac{4}{a^{2}}+\dfrac{2}{b^{2}}=1$．将$b^{2}=a^{2}-4$代入上式可得$\dfrac{4}{a^{2}}+\dfrac{2}{a^{2}-4}=1$．  
去分母化简得$a^{4}-10a^{2}+16=0$，解得$a^{2}=8$或$a^{2}=2$．  

因为$a^{2}>c^{2}=4$，所以$a^{2}=8$，从而$b^{2}=4$．所以该椭圆的标准方程为$\dfrac{x^{2}}{8}+\dfrac{y^{2}}{4}=1$．  

（2）  
设直线$l$的方程为$y=kx+m$，设$A(x_{1},y_{1})$，$B(x_{2},y_{2})$，$M(x_{0},y_{0})$．  

将点$A$与点$B$的坐标分别代入椭圆方程得$\dfrac{x_{1}^{2}}{a^{2}}+\dfrac{y_{1}^{2}}{b^{2}}=1$且$\dfrac{x_{2}^{2}}{a^{2}}+\dfrac{y_{2}^{2}}{b^{2}}=1$，两式相减得$\dfrac{x_{1}^{2}-x_{2}^{2}}{a^{2}}+\dfrac{y_{1}^{2}-y_{2}^{2}}{b^{2}}=0$．利用平方差公式化简得$\dfrac{(x_{1}-x_{2})(x_{1}+x_{2})}{a^{2}}+\dfrac{(y_{1}-y_{2})(y_{1}+y_{2})}{b^{2}}=0$．  

因为直线$l$与椭圆交于不重合的两点，所以$x_{1}\neq x_{2}$，两边同除以$x_{1}-x_{2}$可得  
$$
\dfrac{x_{1}+x_{2}}{a^{2}}+\dfrac{y_{1}-y_{2}}{x_{1}-x_{2}}\cdot\dfrac{y_{1}+y_{2}}{b^{2}}=0.
$$  

根据中点坐标公式可得$x_{1}+x_{2}=2x_{0}$且$y_{1}+y_{2}=2y_{0}$，又因为直线的斜率为$k$，故有$\dfrac{y_{1}-y_{2}}{x_{1}-x_{2}}=k$．代入上式可得$\dfrac{2x_{0}}{a^{2}}+k\cdot\dfrac{2y_{0}}{b^{2}}=0$，化简得$\dfrac{x_{0}}{a^{2}}+\dfrac{ky_{0}}{b^{2}}=0$．  

当$k\neq0$时，动点$M$的坐标满足方程$y=-\dfrac{b^{2}}{a^{2}k}x$，因为$a$与$b$以及$k$均为非零常数，所以$-\dfrac{b^{2}}{a^{2}k}$为常数，故动点$M$在定直线$y=-\dfrac{b^{2}}{a^{2}k}x$上，且该直线经过原点．  

当$k=0$时，直线$l$平行于$x$轴，由椭圆的几何对称性可得其相交线段的中点$M$必定在$y$轴上．此时动点$M$的坐标满足方程$x=0$，该定直线同样经过原点．  

综上所述，动点$M$始终在一条过原点的定直线上．  

（3）  
作图步骤如下．  
第一步，在给定椭圆内利用直尺作两条不重合的平行弦$AB$与$CD$．  
第二步，分别取弦$AB$与弦$CD$的中点$M_{1}$与$M_{2}$．  
第三步，连接$M_{1}M_{2}$并延长作直线．  
第四步，在同一椭圆内变换方向，再作两条不重合的平行弦$EF$与$GH$，且与$AB$不平行．  
第五步，分别取弦$EF$与弦$GH$的中点$M_{3}$与$M_{4}$．  
第六步，连接$M_{3}M_{4}$并延长作直线．  
第七步，直线$M_{1}M_{2}$与直线$M_{3}M_{4}$的交点即为该椭圆的中心$O$．

\begin{tikzpicture}[scale=1.2]
    % 1. 基础参数与椭圆
    \def\a{3.2}     % 半长轴
    \def\b{2.0}     % 半短轴
    \def\rot{30}    % 椭圆旋转角度
    
    % 绘制椭圆并命名路径以便求交点
    \draw[thick, rotate=\rot, name path=ellipse] (0,0) ellipse (\a cm and \b cm);
    
    % 中心 O
    \coordinate (O) at (0,0);
    \fill (O) circle (1.2pt) node[below right] {$O$};

    % 2. 第一组平行弦 AB (上) 和 CD (下)
    \def\angOne{18} % 斜率角度
    % 通过调整基准点让 A 落在椭圆最左侧附近
    \path[name path=lineAB] (0, 0.6) +(\angOne:-5) -- +(\angOne:5);
    \path[name path=lineCD] (0, -1.2) +(\angOne:-5) -- +(\angOne:5);
    
    % 自动求交点
    \path [name intersections={of=ellipse and lineAB, by={A,B}}];
    \path [name intersections={of=ellipse and lineCD, by={C,D}}];
    
    % 绘制弦并打上黑点和标签
    \draw (A) -- (B);
    \draw (C) -- (D);
    \fill (A) circle (1.2pt) node[left]  {$A$};
    \fill (B) circle (1.2pt) node[right] {$B$};
    \fill (C) circle (1.2pt) node[left]  {$C$};
    \fill (D) circle (1.2pt) node[right] {$D$};
    
    % 中点 M1, M2
    \coordinate (M1) at ($(A)!0.5!(B)$);
    \coordinate (M2) at ($(C)!0.5!(D)$);
    \fill (M1) circle (1.2pt) node[above  right] {$M_1$};
    \fill (M2) circle (1.2pt) node[below right] {$M_2$};

    % 3. 第二组平行弦 EF (左) 和 GH (右)
    \def\angTwo{82} % 近乎竖直的倾斜角度
    \path[name path=lineEF] (-0.6, 0) +(\angTwo:-5) -- +(\angTwo:5);
    \path[name path=lineGH] (1.2, 0) +(\angTwo:-5) -- +(\angTwo:5);
    
    % 自动求交点（注意由于角度是从左下到右上，第一个交点是底部的）
    \path [name intersections={of=ellipse and lineEF, by={E,F}}];
    \path [name intersections={of=ellipse and lineGH, by={G,H}}];
    
    % 绘制弦并打上黑点和标签
    \draw (E) -- (F);
    \draw (G) -- (H);
    \fill (E) circle (1.2pt) node[above] {$E$};
    \fill (F) circle (1.2pt) node[below left] {$F$};
    \fill (G) circle (1.2pt) node[above right] {$G$};
    \fill (H) circle (1.2pt) node[below] {$H$};
    
    % 中点 M3, M4
    \coordinate (M3) at ($(E)!0.5!(F)$);
    \coordinate (M4) at ($(G)!0.5!(H)$);
    \fill (M3) circle (1.2pt) node[below left] {$M_3$};
    \fill (M4) circle (1.2pt) node[below right] {$M_4$};

    % 4. 绘制延伸到椭圆边界的中点连线 (完美解决不相交问题)
    % 构造足够长的射线求与椭圆的交点
    \path[name path=rayM1M2] ($(M1)!-2!(M2)$) -- ($(M1)!3!(M2)$);
    \path [name intersections={of=ellipse and rayM1M2, by={M1ext1, M1ext2}}];
    \draw (M1ext1) -- (M1ext2); % 画出贯穿椭圆的直线

    \path[name path=rayM3M4] ($(M3)!-2!(M4)$) -- ($(M3)!3!(M4)$);
    \path [name intersections={of=ellipse and rayM3M4, by={M3ext1, M3ext2}}];
    \draw (M3ext1) -- (M3ext2); % 画出贯穿椭圆的直线

    % 5. 绘制辅助虚线 AC 及其延长线
    % 利用向量计算，让虚线向 A 上方和 C 下方延伸出 15% 的长度
    \coordinate (dirAC) at ($(C)-(A)$);
    \coordinate (A_up) at ($(A) - 0.15*(dirAC)$);
    \coordinate (C_down) at ($(C) + 0.15*(dirAC)$);
    
    \draw[dashed] (A_up) -- (C_down);

    % 6. 精确的角标记 (同旁内角)
    % A 处：弦 AB 与 虚线 AC 的夹角
    % C 处：弦 CD 与 虚线 CA 的夹角
    \pic [draw, angle radius=0.25cm] {angle = B--A--A_up};
    \pic [draw, angle radius=0.35cm] {angle = B--A--A_up};
    \pic [draw, angle radius=0.25cm] {angle = D--C--A};
    \pic [draw, angle radius=0.35cm] {angle = D--C--A};

\end{tikzpicture}
\end{solutions}